\documentclass[11pt,a4paper]{article}
\usepackage[T1]{fontenc}
\usepackage[utf8]{inputenc}
\usepackage[margin=22mm]{geometry}
\usepackage{lmodern}
\usepackage{amsmath,amssymb}
\usepackage{enumitem}
\usepackage{xcolor}
\usepackage{microtype}
\usepackage{hyperref}
\definecolor{epflred}{HTML}{E3000B}
\hypersetup{colorlinks=true,urlcolor=epflred,linkcolor=epflred,
  pdftitle={ENG-654 Project 8: Redundancy and Feasible IK Families in 7R Robots}}
\setlength{\parindent}{0pt}
\setlength{\parskip}{6pt}
\setlist[itemize]{leftmargin=*,itemsep=5pt,topsep=4pt}
\urlstyle{tt}

\begin{document}
{\small\textbf{EPFL \textbar\ ENG-654}\hfill Kinematics-Grounded Motion Planning for Robots}
\vspace{5mm}

{\color{epflred}\Large\bfseries Project 8}

{\LARGE\bfseries Redundancy and Feasible IK Families in 7R Robots\par}

\textbf{Format:} Group of two students; simulation-based project.\\
\textbf{Prerequisites:} Rigid transforms, D--H modelling, redundant inverse
kinematics, Jacobian rank, joint limits, and path continuation.

\section*{Motivation}
For a redundant robot, a fixed tool pose can be realized by a continuous
family of postures. A selected redundant angle provides a useful way to inspect
that family, but joint limits and singularities restrict it. Understanding
those restrictions is necessary before choosing a redundancy schedule along
a workspace path.

\section*{Task}
\textbf{Overall objectives.} Analyze a small family of spherical-shoulder/elbow/wrist 7R manipulators
and explain how geometry changes fixed-pose IK families, their classification,
and feasibility over a short path.

\begin{itemize}
  \item \textbf{Define a bounded geometric family.} Use iiwa-like geometry as a
reference and vary only the upper-arm to forearm length ratio over three
instructor-supplied parameter sets with the same total arm length. The
parameterized FK/IK implementation will be provided. Verify the reference
against its URDF and describe the synthetic variants using their actual
kinematic frames rather than unchanged meshes.

  \item \textbf{Inspect fixed-pose self-motion.} At two selected full poses, sample
$\lambda=q_3$ and evaluate all returned IK candidates. Plot joint values versus
$\lambda$, check full-pose closure, and identify intervals removed by limits
or loss of regularity. Explain the locally one-dimensional nature of the
regular fixed-pose solution set.

  \item \textbf{Classify the recovered families.} Preserve the supplied analytical
branch labels, document their geometric meaning where available, and track
their continuation across neighboring samples. Distinguish a label within a
fixed-angle slice from a connected component of the complete self-motion set.
Do not use solver array order as identity.

  \item \textbf{Generate comparative feasibility charts.} For one common short pose
path, plot path sample horizontally and redundancy angle vertically, using a
separate panel for each labelled candidate family. Distinguish absent IK,
pose error, joint-limit violation, full-Jacobian rank loss, and reduced-chart
failure. Compare the three geometries on identical axes.

  \item \textbf{Interpret feasible connections.} Extract one candidate corridor or
explain the sampled obstruction for each geometry. Re-solve intermediate poses
before accepting neighboring chart cells as connected. Compare feasible-angle
width and minimum joint-limit/regularity margins, and explain what the chosen
chart can miss.
\end{itemize}

\newpage
\section*{Specifics and scope}
Use three geometries, two fixed poses, and one short path, initially with
about 100 pose samples and 120 redundancy angles. Refine one apparent boundary
or connection. This is a geometric family study, not a full trajectory-optimization
project. The provided solver must support all three parameter sets; do not
apply an unmodified iiwa-specific solver to a changed geometry.

For every chart, keep pose validity, joint limits,
full-Jacobian regularity, and reduced-chart regularity as separate diagnostics.
The rectangular $6\times7$ Jacobian has no determinant. A feasible pixel is
not a validated edge: re-solve intermediate prescribed poses and check joint
continuity within limits. If collision checking is unavailable, state that
feasibility is kinematic only; a clearance proxy is not a full collision test.

\section*{Expected outcome}
Understand the URDF reference, the parameterized D--H geometry, redundant
IK families, and the relation between classification and feasibility. Deliver
self-motion plots and branch-labelled path-sample versus redundant-angle charts
that explain how the selected geometric change affects feasible connections.

Submit reproducible code, model parameters and test data, the required plots,
a concise 5--6-page report, and a short simulation demonstration. Explain
both successful cases and failures; conclusions must follow the numerical
and geometric evidence rather than an expected outcome.

\section*{Provided materials and implementation}
The KUKA LBR iiwa 7 robot description (.urdf) and link meshes (.stl) will be
provided to the students. An additional URDF with unrestricted revolute joints
will be provided for mathematical continuation; apply physical limits from
the reference model separately. The three geometric parameter sets, compatible
parameterized analytical FK/IK implementation, and reference data will also
be supplied. Reuse the IK implementation rather than deriving a general 7R
solver. The reference data validate the original iiwa member only; check
the synthetic variants against their own forward models.

Construct each chart by evaluating IK at every sampled pose and redundancy
angle, applying the stated validity and regularity checks, and tracking the
provided labels across neighboring samples. Use a robotics library or your own
tools to inspect frames, verify D--H parameters, visualize self-motion, and
animate candidate paths.

\end{document}
